%% 毕业设计小结

\begin{designsummary}
本毕业设计以“基于深度学习的 B-rep 端到端生成算法研究”为题，目标是在不依赖显式 VQ-VAE 码本的前提下，验证一种点级 NURBS 序列化 + 标准 next-token 自回归 Transformer 的端到端 B-rep 生成框架。围绕这一目标，整个工作可以归纳为研究、实现、实验与写作四个阶段。

研究阶段，通过对 PolyGen、DeepCAD、SkexGen、BrepGen 等代表性工作的系统阅读，梳理出 CAD 操作序列、草图—拉伸解耦码本、结构化潜空间扩散三条主线在中间表征上各自的优势与限制，进而把本课题的研究定位为“点级量化 + 自回归 Transformer + OpenCASCADE 逆向重建”的工程化探索。本阶段同步完成了 B-rep 拓扑层级、NURBS 数学定义、Transformer / LLaMA3 关键组件、OpenCASCADE BRepCheck 评测口径等技术铺垫。

实现阶段，完成了数据预处理（STEP $\to$ 拓扑分解 $\to$ NURBS 采样 $\to$ 量化 $\to$ 拓扑感知排序）、自回归生成模型（LLaMA3 风格主干 + 标准 next-token 头，CompoundPointDecoder 作为可选扩展保留）、训练流程（PyTorch Lightning + DDP + AMP）与推理（top-k / top-p 自回归采样）、逆向重建（NURBS 拟合 + 拓扑缝合 + BRepCheck 校验）等全部环节的代码实现。数据词表规模 4\,053（50 面索引 + 4\,000 量化坐标 + 3 个结构 token），模型 \texttt{vocab\_size = 4\,054}（额外引入 PAD 用于批内对齐），face\_block = 49（16 个三维点 + 1 个面索引），edge\_block = 14（2 个面索引 + 4 个三维点）。

实验阶段，在 DeepCAD 子集（训练 64\,351 / 验证 6\,096）上完整训练 200 个 epoch（4 张 NVIDIA H20，端到端 17 小时 28 分钟）。最终模型 hr2\_patch1 验证集交叉熵损失收敛至 0.10486，1\,005 个独立生成样本经反向重建后 BRepCheck 通过率达 97.51\%（of saved），早期 hr2 基线在 3\,000 个写出样本上的对应数字为 90.40\%，两者均显著超过毕业设计任务书 60\% 的红线指标。同时记录了未产出 / 产出但无效 / 产出且有效三种状态的样本分布，并明确指出本文未单独对排序策略与 CompoundPointDecoder 等组件做受控对照消融，因此不对其独立贡献做定量声明。

写作阶段，按照“研究问题 $\to$ 相关技术 $\to$ 序列化与预处理 $\to$ 自回归模型与重建 $\to$ 实验与总结”五章结构组织全文，并严格遵守“每一个关键技术主张必须可追溯到文献、代码、实验日志或学校任务书”的证据规范，对于尚未在本论文中完成的实验（如跨数据集泛化、排序受控消融、CompoundPointDecoder 启用对照、几何相似度评估、VQ-VAE 路线复现等）一律放入论文末尾的局限性与后续工作部分，宁缺勿造。

通过本毕业设计，本人在 CAD 几何表征、深度生成模型工程实现、大规模训练实验组织、学术写作与证据管理等方面均获得了系统性的训练；同时也清楚地认识到本课题在跨数据集泛化、各组件独立贡献量化、几何相似度评估指标等方向上仍有进一步深入的空间，相关后续工作已在第~5 章末尾以条目形式列出。
\end{designsummary}
